% Methods (main text, after the Discussion): a summary. The full protocol is in the Supplementary Methods (supplement.tex).
\section{Methods}
\label{sec:setup}\label{sec:mmethods}

Every number, table and figure in this paper is produced by one pipeline from public data
(\texttt{papers/manuscripts/s1-format/bench} in the \texttt{singlet} repository); the Supplementary Methods give the full
protocol.

\para{Data}\nm{nMatrices} count matrices from \nm{nDatasets} public datasets in \nm{nProtocolClasses} protocol classes
(droplet, plate-based, combinatorial, multimodal and spatial data, from a thousand to more than a million cells;
Supplementary Table~\ref{S-tab:datasets}) were downloaded from their primary repositories, checked against recorded
SHA-256 digests and converted to one canonical sparse matrix with the source's gene identifiers and barcodes. Matrices with
non-integer values were rejected rather than rounded.

\para{Formats}Each format was written from the canonical matrix by its own library with default or recommended settings
(Supplementary Table~\ref{S-tab:formats}), read back and compared with the canonical matrix. File size includes gene
identifiers and barcodes.

\para{Read speed}Each read ran in a fresh process pinned to \nm{threadsHeadline} cores of a cluster node (\nm{cpu}), after
the library had been loaded and used once, untimed; files were evicted from the operating system's cache before each of
three timed repeats, and we report medians. The core sweep read the same files on 1 to 44 cores of an otherwise idle
48-core node.

\para{Training loaders}Loaders read minibatches of 512 cells, each made of eight blocks of 64 consecutive cells at random
offsets, into PyTorch sparse tensors, from a local S3-compatible object store whose bandwidth (0.1 to 10\,Gbit/s) and added
latency (0 to 200\,ms) we controlled, and from Cloudflare R2 over the public internet. We recorded the time from a new
process to the first minibatch and the time of each later one, for three seeds. Costs per epoch use the bytes and requests
counted at the emulated store and providers' published prices.

\para{Size floor}Poisson, negative binomial and nonparametric count models conditioned on cell depth, and on cell cluster,
were fitted on half of each matrix's cells and scored on the other half, zeros
included~\cite{hafemeister2019normalization,townes2019feature}. The shortest code length among them, charged nothing for
the models' own parameters, is our reference for the floor; it is a reference, not a bound.

\para{Speed frontier}A C++ benchmark encoded each matrix's row-index steps and values with raw 32-bit arrays, fixed-width
bytes, LEB128 varints, Stream VByte~\cite{lemire2018stream}, SIMD bit-packing~\cite{lemire2015decoding}, an order-0 rANS
coder~\cite{duda2013asymmetric}, and zstd at levels 3 and 19, and decoded each, and the \sone{} member, on one pinned core
(AVX2), checked against the input. Only writing the output arrays gives the rate no decoder can exceed by much; the
\emph{ideal code} stores the floor's code length and decodes at that rate. The trend is a power law, fitted to the other
codes' time per entry against their size on each matrix's frontier and evaluated at \sone's size, with an interval from
1,000 bootstrap resamples of the matrices.

\para{Write cost}Each format was written from the canonical matrix in memory, in a fresh process, recording wall time and
peak resident memory beyond the loaded input, each relative to writing uncompressed CSC arrays of the same matrix.
